\documentclass[a4paper,11pt]{article}
\usepackage[utf8]{inputenc}
\usepackage{amsmath}
\usepackage{amssymb}
\usepackage{geometry}
\usepackage{enumitem}
\usepackage{xcolor}

\geometry{top=2.5cm, bottom=2.5cm, left=2.5cm, right=2.5cm}

\title{\textbf{Advanced Quantum Mechanics - Review Mock Exam}}
\author{Based on Lecture 15 Review}
\date{Winter Semester 2025/2026}

\begin{document}

\maketitle

\section*{Instructions}
\begin{itemize}
    \item This exam focuses on the theoretical concepts from Lectures 1 to 14.
    \item Each question has \textbf{5 options}.
    \item \textbf{Any number of options (from 0 to 5)} may be correct for a given question.
    \item Select all statements that are physically and mathematically correct.
\end{itemize}

\hrule
\vspace{1cm}

% Question 1: Fundamentals
\noindent\textbf{1. The Schrödinger Equation and Time Evolution}\\
Consider the general formalism of quantum mechanics as reviewed in the slides. Which of the following statements are correct?

\begin{enumerate}[label=\Alph*.]
    \item The time-evolution operator $\hat{U}(t) = e^{-\frac{i}{\hbar}\hat{H}t}$ is unitary if and only if the Hamiltonian $\hat{H}$ is Hermitian.
    \item For a time-independent Hamiltonian, the probability density $|\psi(x,t)|^2$ of a stationary state $|E_n\rangle$ changes sinusoidally with time.
    \item The general solution to the time-dependent Schrödinger equation can be expanded as a superposition of energy eigenstates: $|\psi(t)\rangle = \sum_n c_n(0) e^{-iE_n t/\hbar} |E_n\rangle$.
    \item The norm of the state vector $\langle \psi(t) | \psi(t) \rangle$ is conserved during unitary time evolution.
    \item In the position representation, the momentum operator acts as a multiplication operator ($\hat{p} \to p$).
\end{enumerate}
\textit{Correct Answers: A, C, D}

\vspace{0.5cm}

% Question 2: Harmonic Oscillator
\noindent\textbf{2. The Quantum Harmonic Oscillator}\\
Regarding the 1D quantum harmonic oscillator and its algebraic solution using ladder operators $\hat{a}$ and $\hat{a}^\dagger$:

\begin{enumerate}[label=\Alph*.]
    \item The energy spectrum is equidistant with $E_n = \hbar\omega(n + 1/2)$.
    \item The ground state energy is zero ($E_0 = 0$).
    \item The creation operator $\hat{a}^\dagger$ increases the particle number/excitation level by 1, such that $\hat{a}^\dagger |n\rangle = \sqrt{n+1} |n+1\rangle$.
    \item The position operator $\hat{x}$ can be written as a linear combination of $\hat{a}$ and $\hat{a}^\dagger$.
    \item The ground state wavefunction $\langle x | 0 \rangle$ is an odd function of $x$ (antisymmetric).
\end{enumerate}
\textit{Correct Answers: A, C, D}

\vspace{0.5cm}

% Question 3: Hydrogen Atom
\noindent\textbf{3. The Hydrogen Atom}\\
Consider the solution to the Hydrogen atom problem ($V(r) \propto -1/r$):

\begin{enumerate}[label=\Alph*.]
    \item The energy eigenvalues $E_n$ depend only on the principal quantum number $n$ (ignoring fine structure).
    \item The degeneracy of the energy level $E_n$ is $n^2$ (including spin degeneracy, it would be $2n^2$, but neglecting spin it is $n^2$).
    \item The wavefunctions are factorized into a radial part $R_{nl}(r)$ and spherical harmonics $Y_{lm}(\theta, \phi)$.
    \item The $l=0$ (s-orbitals) have non-zero probability density at the nucleus ($r=0$).
    \item The centrifugal barrier term in the radial equation vanishes for $l=0$.
\end{enumerate}
\textit{Correct Answers: A, B, C, D, E}

\vspace{0.5cm}

% Question 4: Perturbation Theory
\noindent\textbf{4. Approximation Methods (Perturbation Theory)}\\
Which of the following statements regarding time-independent and time-dependent perturbation theory are true?

\begin{enumerate}[label=\Alph*.]
    \item The first-order energy correction in time-independent perturbation theory is the expectation value of the perturbation in the unperturbed state: $E_n^{(1)} = \langle n^{(0)} | \hat{V} | n^{(0)} \rangle$.
    \item Fermi's Golden Rule gives the transition rate for a system interacting with a continuum of final states under a constant or harmonic perturbation.
    \item Second-order perturbation theory always leads to an increase in the ground state energy.
    \item To apply non-degenerate perturbation theory, the energy difference between states must be large compared to the matrix element of the perturbation.
    \item In the interaction picture, the state vector is constant in time.
\end{enumerate}
\textit{Correct Answers: A, B, D}

\vspace{0.5cm}

% Question 5: Scattering Theory
\noindent\textbf{5. Scattering Theory and Partial Waves}\\
Regarding scattering by a spherically symmetric potential:

\begin{enumerate}[label=\Alph*.]
    \item The First Born approximation is generally most accurate at very low incident energies ($k \to 0$).
    \item In the low-energy limit ($kR \ll 1$), s-wave ($l=0$) scattering usually dominates the cross-section.
    \item The differential cross-section is related to the scattering amplitude by $d\sigma/d\Omega = |f(\theta)|^2$.
    \item Phase shifts $\delta_l(k)$ are necessary to calculate the scattering amplitude in the partial wave method.
    \item The scattering amplitude $f(\theta)$ is purely real for all potentials.
\end{enumerate}
\textit{Correct Answers: B, C, D}

\vspace{0.5cm}

% Question 6: Identical Particles
\noindent\textbf{6. Many-Body Physics and Identical Particles}\\
Consider a system of $N$ identical particles:

\begin{enumerate}[label=\Alph*.]
    \item For bosons, the total wavefunction must be antisymmetric under the exchange of any two particles.
    \item Fermions obey the Pauli exclusion principle, meaning no two fermions can occupy the same single-particle state.
    \item The ground state of two non-interacting bosons in a harmonic trap has energy $E_{gr} = \hbar\omega$.
    \item The exchange symmetry requirement arises because identical particles are indistinguishable in quantum mechanics.
    \item The Slater determinant is used to construct symmetric wavefunctions for bosons.
\end{enumerate}
\textit{Correct Answers: B, C, D}

\vspace{0.5cm}

% Question 7: Molecules and Atoms
\noindent\textbf{7. H$_2^+$ and Helium Atom}\\
Regarding the application of quantum mechanics to simple molecules and multi-electron atoms:

\begin{enumerate}[label=\Alph*.]
    \item The Born-Oppenheimer approximation relies on the fact that the nuclei are much heavier (and thus move slower) than the electrons.
    \item In the H$_2^+$ molecule, the bonding orbital is antisymmetric with respect to the inversion of coordinates.
    \item The ground state of the Helium atom is a spin singlet state (para-helium).
    \item The Hartree-Fock method captures all electron-electron correlation effects perfectly.
    \item The variational method states that any trial wavefunction yields an energy expectation value greater than or equal to the true ground state energy.
\end{enumerate}
\textit{Correct Answers: A, C, E}

\vspace{0.5cm}

% Question 8: Second Quantization
\noindent\textbf{8. Second Quantization formalism}\\
In the formalism of second quantization:

\begin{enumerate}[label=\Alph*.]
    \item It is a necessary framework to describe systems where the number of particles can change.
    \item Bosonic creation and annihilation operators satisfy commutation relations $[\hat{a}_i, \hat{a}_j^\dagger] = \delta_{ij}$.
    \item Fermionic field operators satisfy anticommutation relations $\{\hat{\Psi}(\mathbf{x}), \hat{\Psi}^\dagger(\mathbf{y})\} = \delta(\mathbf{x}-\mathbf{y})$.
    \item The operator $\hat{N} = \sum_i \hat{a}_i^\dagger \hat{a}_i$ counts the total number of particles in the system.
    \item A two-body interaction term involves four operators (e.g., $\hat{a}^\dagger \hat{a}^\dagger \hat{a} \hat{a}$).
\end{enumerate}
\textit{Correct Answers: A, B, C, D, E}

\end{document}